\section*{第 7 章：无类型 lambda 演算的 ML 实现}
\addcontentsline{toc}{section}{第 7 章：无类型 lambda 演算的 ML 实现}

\begin{taplsolutionentrybox}
\phantomsection
\label{sol:translator:7.3.1}
\textbf{练习 7.3.1}

按值调用的大步求值器先把函数位置求值为抽象，再把参数求值为值，执行一次
顶层替换，最后递归求值替换结果：

\begin{taplocamlcode}
let rec eval ctx t =
  match t with
    TmAbs _ -> t
  | TmApp(_, t1, t2) ->
      let v1 = eval ctx t1 in
      let v2 = eval ctx t2 in
      (match v1 with
         TmAbs(_, _, t12) ->
           eval ctx (termSubstTop v2 t12)
       | _ -> raise NoRuleApplies)
  | _ -> raise NoRuleApplies
\end{taplocamlcode}
\taplrustcounterpart{sol-translator-07-eval-big}

对闭项而言，这段程序与练习 5.3.8 的规则逐项对应：\texttt{TmAbs} 对应
\textsc{B-Value}，\texttt{TmApp} 分支对应 \textsc{B-App} 的三个前提。
若对含自由变量的开项求值，变量没有大步规则，因而抛出
\texttt{NoRuleApplies}。

\taplsolutionbacklink{ex:7.3.1}{返回练习 7.3.1}
\end{taplsolutionentrybox}
